\begin{definition}
Дерево — непустой связный граф без циклов. В этом билете граф простой.
\(\chi_G(k)\) — число правильных раскрасок различимыми цветами из \([k]\),
продолженное до многочлена по билету 18.
\end{definition}
\begin{lemma}[Удаление листа]
Если \(v\) — лист, то \(\chi_G(k)=(k-1)\chi_{G-v}(k)\).
\end{lemma}
\begin{proof}
Для каждой раскраски \(G-v\) цвет единственного соседа \(v\) запрещён,
а остальные \(k-1\) цветов разрешены. Продолжения независимы от других вершин.
Равенство чисел раскрасок для всех натуральных \(k\) даёт тождество многочленов.
\end{proof}
\begin{theorem}[Характеристика дерева]
Для простого графа на \(n\ge1\) вершинах
\[
 \boxed{G\text{ — дерево}\quad\Longleftrightarrow\quad
 \chi_G(k)=k(k-1)^{n-1}.}
\]
\end{theorem}
\begin{proof}
Необходимость: при \(n=1\) раскрасок \(k\). При \(n>1\) удаляем лист;
оставшийся граф — дерево. Индукция и лемма дают
\((k-1)\cdot k(k-1)^{n-2}=k(k-1)^{n-1}\).
Можно также выбрать корень: его цвет произволен, а каждая другая
вершина выбирает любой цвет, кроме цвета родителя. В дереве других
ограничений между уже выбранными предками нет.

Достаточность: степень многочлена равна числу вершин по билету 18.
Коэффициент при \(k\) у \(k(k-1)^{n-1}\) равен \((-1)^{n-1}\ne0\).
По лекционной теореме билета 20 наименьшая ненулевая степень равна
числу компонент (её доказательство не использует данный билет).
Значит, \(G\) связен.
Для простого графа коэффициент при \(k^{n-1}\) равен \(-m\)
(доказательство дано в билете 20 и не использует данный билет).
Сравнение с разложением \(k(k-1)^{n-1}\) даёт \(m=n-1\)
при \(n\ge2\). Связный граф с \(n-1\) рёбрами — дерево по билету 6.
При \(n=1\) простой граф сам является одноточечным деревом.
\end{proof}
\begin{cor}[Многочлен леса]
Для леса с \(n\) вершинами и \(c\) компонентами
\[
 \chi_G(k)=k^c(k-1)^{n-c}.
\]
\end{cor}
\begin{proof}
Каждая компонента — дерево. Перемножаем её формулу;
сумма чисел вершин равна \(n\), число множителей \(k\) равно \(c\).
Для пустого леса получаем \(1\).
\end{proof}
\textbf{Зачем нужна простота.} Если разрешить кратные рёбра, удвоение
любого ребра дерева не меняет хроматический многочлен, но полученный
мультиграф уже не дерево в обычном смысле. Для мультиграфов формула
характеризует дерево, полученное после удаления кратностей, при отсутствии петель.
